%% ma: L12-B14-0098
%% lop: 12
%% chuong: 5
%% bai: 14
%% dang: 12.14.28
%% loai: DS
%% muc: [TH, TH, TH, VD]
%% dap_an: "ĐSĐS"
%% nguon: "Ảnh giáo viên gửi 07/11/2026 (Tư Duy Mở - tuduymo.com, đề tự luyện Phương trình mặt phẳng, đề lập phương trình mặt phẳng), Câu 15"
%% kien_thuc: "Bài 14 (phương trình mặt phẳng tổng quát, giao điểm với trục tọa độ)"
%% trang_thai: chinh-thuc
%% ghi_chu: "Tài liệu gốc không kèm đáp án; đáp án do Claude giải và kiểm bằng sympy. Dạng đã duyệt 07/11/2026."
\begin{ex}
Trong không gian $Oxyz$, cho điểm $A(a;2b-3;10-a-b)$, với $a$ và $b$ là những tham số thực thay đổi. Khi $a$, $b$ thay đổi thì $A$ luôn thuộc một mặt phẳng $(\alpha)$ cố định.
\choiceTF
{\True Giao điểm của $(\alpha)$ với trục $Oy$ có tọa độ là $(0;17;0)$.}
{Giao điểm của $(\alpha)$ với trục $Ox$ có tọa độ là $(2;0;0)$.}
{\True Phương trình tổng quát của mặt phẳng $(\alpha)$ là $2x+y+2z-17=0$.}
{Có tất cả $8$ điểm có các thành phần tọa độ là các số tự nhiên và thuộc $(\alpha)\cap(Oxy)$.}
\loigiai{$x=a$, $y=2b-3\Rightarrow b=\dfrac{y+3}2$, $z=10-x-\dfrac{y+3}2\Rightarrow(\alpha)\colon 2x+y+2z-17=0$.

a) $x=z=0\Rightarrow y=17$. Đúng.

b) $y=z=0\Rightarrow x=\dfrac{17}2$. Sai.

c) Đúng.

d) $z=0$, $2x+y=17$ với $x,y\in\mathbb N$: $x\in\{0;1;\ldots;8\}$, có $9$ điểm. Sai.}
\end{ex}
